%% kapa: exam=1 round=21 year=2010 period=1 form=A subject=civil_law subject_ko=민법 seq=6 number=6
%% source: Q-Net 감정평가사 기출문제 (cst00309) artlSeq=5024181 file=extracted/1008/5024181/2010년 제21회 감정평가사 1차시험 기출문제/1교시.hwp
%% conversion: kordoc 4.24.51 -> tools/gen1.py (auto); review=proofread-sonnet; flags=none
\begin{question}{6}
등기의 추정력에 관한 설명으로 옳지 않은 것은? (다툼이 있으면 판례에 의함)
\begin{choices}
  \item 소유권이전등기가 경료되어 있는 경우, 그 등기명의자는 제3자 뿐만 아니라 전(前) 소유자에 대하여도 적법한 등기원인에 의하여 소유권을 취득한 것으로 추정된다.
  \item 저당권의 등기는 그 저당권의 존재 뿐만 아니라 피담보채권의 존재도 추정하게 한다.
  \item 부동산의 표시에 관한 사항에는 추정력이 인정되지 않는다.
  \item 등기의무자가 사망하기 전에 등기원인이 이미 존재하였다면 등기의무자의 사망 후에 그로부터 경료된 등기라도 적법한 것으로 추정된다.
  \item 소유권이전등기의 원인으로 주장된 계약서가 진정하지 않은 것으로 증명된 경우에도 다른 적법한 등기원인이 있을 것으로 추정된다.
\end{choices}
\end{question}
\answer{5}
